\cvheader

I study how people, building systems, and climate interact to shape thermal protection and
energy demand. My work asks how building decisions can account for the likelihood and consequences
of unacceptable discomfort across diverse occupants and uncertain climate futures.

\section{Research programme}
\label{sec:research-programme}
{\fontsize{9}{11}\selectfont\color{muted}%
\ifdefined\cvShortProfile
  Selected work links to the full CV bibliography.
\else
  Selected work links to the bibliography below.
\fi\par}
\thread{People / Thermal protection and human variability}
  {Predicting how hot or cold people feel, accounting for human heat generation and context.}
  {\pubref{J10}{human heat defaults}; \pubref{J07}{contextual drivers};
   \pubref{J06}{thermal sensation probabilities}.}
\thread{Systems / Integration and energy uncertainty}
  {Linking human response and building control to assess discomfort risk alongside energy demand.}
  {\pubref{J09}{co-simulation}; \pubref{A02}{model benchmarking};
   \pubref{A01}{HVAC supervision using sensation probabilities}.}
\thread{Climate / Future weather and resilience}
  {Future weather, input fidelity, and uncertainty in climate-resilient building performance.}
  {sAMY --- \pubref{J03}{generation} and \pubref{J05}{input fidelity};
   \pubref{A03}{cooling and ageing}.}

\section{Appointments}
\cvrow{2023--present}{Assistant Professor / The University of Hong Kong}
  {Department of Architecture; appointed September 2023.}
\cvrow{2020--2023}{Principal Data Scientist / Bank of New York Mellon}
  {Led ML deployment and model governance for daily USD 29--43 billion Federal Reserve account-balance forecasts.}
\cvrow{2019--2020}{Postdoctoral Researcher / Princeton University}
  {Building science, sensing, and environmental modeling.}

\section{Education}
\cvrow{2019}{Ph.D., Architecture / Princeton University}
  {Architectural Technology \& Computational Design.}
\cvrow{2013}{M.Eng., Mechanical Engineering / Columbia University}{}
\cvrow{2012}{B.Eng., Architectural Engineering / Harbin Institute of Technology}
  {HVAC Engineering.}

\section{Research funding}
\cvrow{2026--}{ASHRAE Research Project 1959-TRP / \$55,000}
  {Principal investigator; demographic-aware metabolic-rate defaults for thermal comfort standards.
   Awarded February 2026; 15 months from April 2026. Sponsor: TC 4.1.}

\clearpage
\section{Teaching innovation / funding}
\cvrow{2025--2028}{Socratic Oracle / HK\$1,000,000}
  {Principal investigator, Teaching Development \& Language Enhancement Grant, HKU.
   Multimodal voice-AI tutor for evidence-based learning.}
\cvrow{2026--2028}{AI Design Coach / HK\$300,000}
  {Principal investigator, Teaching Development Grant, HKU. Conversational support for design narratives across Architecture, Geography, and Computer Science.}

\section{Honors \& awards}
\cvrow{May 2026}{Best Paper Award / IAQVEC 2026, Los Angeles}
  {Co-recipient with PhD student Kanxuan He, as supervisor and co-author, for
   \href{https://doi.org/10.1051/e3sconf/202671610011}{climate-zone classification for building codes}.}
\cvrow{2025}{1st Place / ICML CO-BUILD Smart Building Competition}
  {Faculty mentor and author.}
\cvrow{2025}{3rd Place / NeurIPS Urban AI Workshop, Buildings Challenge}
  {Faculty mentor.}
\cvrow{2021}{Gartner Eye on Innovation Award / BNY Mellon}{}
\cvrow{2018}{Lowry Methodology Award / International Conference of Urban Climate}{}

\section{Supervision}
\cvrow{2025--present}{Kanxuan He / PhD, HKU}
  {Probabilistic control and climate resilience.}
\cvrow{2024--present}{Yu Chang / PhD, HKU}
  {Cross-cultural thermal comfort; first-author review in
   \emph{Renewable and Sustainable Energy Reviews} (2025).}
\cvrow{2026}{MArch thesis supervision / HKU}
  {Physics-informed occupant modeling; EEG-based neuroarchitecture.}

\section{Teaching}
\cvrow{2025--present}{\href{https://clementinec.github.io/ARCH7476}{ARCH7476 / Evidence-Based Generative Design}}
  {Designed and taught; MArch / UG / PhD. Reproducible computational design.}
\cvrow{2024--present}{\href{https://clementinec.github.io/ccgl9065}{CCGL9065 / Our Response to Climate Change}}
  {Designed and taught; undergraduate. Climate action and public exhibition \emph{Hong Kong 2100}.}
\cvrow{2023--present}{\href{https://clementinec.github.io/desn2003}{DESN2003 / Research for Innovation}}
  {Designed and taught; undergraduate. Research methods and communication.}
\cvrow{2015--2018}{Prior teaching / Princeton University}
  {Mini-courses on IoT sensing and thermodynamics laboratories.}

\section{Service \& public scholarship}
\textbf{Academic leadership.} Co-Chair, Technical Development \& Proceedings Integration,
CAAD Futures 2025.\par
\textbf{Peer review.} Journals including \emph{Energy and Buildings},
\emph{Building and Environment}, and \emph{Applied Energy}; workshops at CVPR, ICML, and NeurIPS.\par
\textbf{Exhibition.} Co-author and exhibitor, \emph{Social Condenser Extraordinaire:
Hong Kong's Municipal Services Buildings}, \emph{Projecting Future Heritage:
A Hong Kong Archive}; Hong Kong collateral exhibition, 19th Venice Architecture
Biennale, May--November 2025.
